\section{User experience and error reporting} \label{sec:user-experience}

We discuss the Verus user experience by example. Suppose the user starts by defining an \verb`Account` struct and
an \verb`exec` function to transfer funds between accounts.

\begin{lstlisting}[language=Verus,style=VerusLineNos,firstnumber=6]
pub struct Account { pub balance: u64 }

pub fn transfer_funds(orig: &mut Account, dest: &mut Account, amount: u64) {
  requires([ old(orig).balance >= amount, old(dest).balance as nat + amount < u64::MAX ]);
  ensures([ dest.balance == old(dest).balance + amount, orig.balance == old(orig).balance - amount ]);
  orig.balance = orig.balance - amount;
  dest.balance = dest.balance + amount;
}
\end{lstlisting}

This function verifies, because Rust's type system ensures that \verb`orig` and \verb`dest`
are not aliased. In fact, if the user accidentally aliased the two arguments when calling \verb`transfer_funds`,
\begin{lstlisting}[language=Verus,style=VerusLineNos,firstnumber=15]
fn main() {
  let mut acct1 = Account { balance: 20_000 };
  transfer_funds(&mut acct1, &mut acct1, 20_000);
  assert(acct1.balance == 10_000);
}
\end{lstlisting}
\noindent the user would quickly get an error from the Rust borrow checker, and Verus would not
attempt to invoke Z3 to verify the invalid program, thereby allowing the user to quickly iterate by fixing
the issue and re-running Verus.
\begin{lstlisting}[style=Verus,literate=]
error[E0499]: cannot borrow `acct1` as mutable more than once at a time
  --> account.rs:17:32
   |
17 |     transfer_funds(&mut acct1, &mut acct1, 20_000);
   |     -------------- ----------  ^^^^^^^^^^ second mutable borrow occurs here
   |     |              |
   |     |              first mutable borrow occurs here
   |     first borrow later used by call
\end{lstlisting}

If one wrote similar code in Dafny, using a \verb`class` (a reference type) to
represent the \verb`Account`, they would declare the \verb`transfer_funds` method as
\lstinline[style=Verus]|method TransferFunds(orig: Accnt, dest: Accnt, amnt: nat)| with
similar preconditions and postconditions. Dafny would report that the postconditions
cannot be verified, which can be misleading to the developer, who has to determine that such
a failure is due to potential aliasing of \verb`orig` and \verb`dest`; the developer
would then need to add a framing condition to \verb`transfer_funds`,
\verb`requires orig != dest`.

In response to the Rust borrow checker failure above, the user may try and fix the \verb`main` function,

\begin{lstlisting}[language=Verus,style=VerusLineNos,firstnumber=15]
fn main() {
  let mut acct1 = Account { balance: 10_000 }; let mut acct2 = Account { balance: 20_000 };
  #[spec] let total_balance = acct1.balance + acct2.balance;
  transfer_funds(&mut acct1, &mut acct2, 20_000);
  assert(total_balance == acct1.balance + acct2.balance);
}
\end{lstlisting}

\noindent but inadvertently introduce a logical error, which Verus reports with precise pointers to the
offending code, and the relevant context (in this case, the failing precondition on the definition
of \verb`transfer_funds`):

\begin{lstlisting}[style=Verus,literate=]
error: precondition not satisfied
  --> account.rs:18:5
   |
9  |   requires([ old(orig).balance >= amount, old(dest).balance as nat + amount < u64::MAX ]);
   |              --------------------------- failed precondition
...
18 |     transfer_funds(&mut acct1, &mut acct2, 20_000);
   |     ^^^^^^^^^^^^^^^^^^^^^^^^^^^^^^^^^^^^^^^^^^^^^^
\end{lstlisting}

Adjusting the transferred amount to \verb`10_000` would result in successful verification upon re-running Verus.

This user experience is conceptually similar to that of the Viper separation logic engine~\cite{DBLP:conf/vmcai/0001SS16}
and the VeriFast C and Java verification tool~\cite{VeriFast}, with the distinction that both of these
tools use a separate substructural logic to reason about memory permissions:
the user has to explicitly manipulate permissions and separation logic predicates for the program data.
In Rust's linear type system, memory-reasoning permissions are implicitly associated with data ownership,
and manipulated by moving values, or taking references.
